\title{Byte Latent Transformer: Patches Scale Better Than Tokens}

\begin{abstract}
We present the Byte Latent Transformer (\textsc{BLT}), an innovative byte-level large language model (LLM) architecture that, for the first time, achieves performance on par with traditional tokenization-based LLMs at scale, while offering marked gains in inference speed and model robustness. \textsc{BLT} operates by grouping bytes into adaptively sized patches, which constitute the primary computational units within the model. The patch boundaries are determined by the entropy of upcoming bytes, allocating greater computational resources and capacity in regions with higher information complexity. We provide the inaugural flop-controlled scaling analysis of byte-level architectures, spanning up to 8 billion parameters and 4 trillion bytes of training data. Our findings confirm that it is viable to train models directly on raw byte sequences, without the need for predetermined vocabularies, and achieve efficient scaling. Both the training process and inference benefit from selecting longer patches when the underlying data exhibits high predictability, and the approach brings qualitative gains in both complex reasoning and long-tail generalization. In summary, for a given inference budget, \textsc{BLT} surpasses the scaling efficiency of token-based models by jointly expanding both patch dimensions and overall model capacity.
\end{abstract}

\section{Introduction}
\label{section:intro}

We present the Byte Latent Transformer~(\textbf{BLT}), a novel architecture that eliminates the need for tokenization by directly learning from raw byte sequences. BLT achieves parity with conventional tokenization-based approaches at scale, while offering marked gains in both efficiency and resilience (\S\ref{sec:robustness}).
Although large language models (LLMs) are generally trained in a predominantly end-to-end manner, tokenization remains a heuristic pre-processing stage wherein bytes are segmented into a predetermined, static vocabulary. This process introduces biases in the compression of input sequences, giving rise to notable issues including modality and domain sensitivity~\citep{dagangetting}, vulnerability to input perturbations~(\S\ref{sec:robustness}), limited orthographic representation~\citep{edman2024cute}, and cross-lingual inequities~\citep{liang2023xlm,petrov2024language,limisiewicz2024myte}.

Historically, tokenization has played a crucial role since training large language models on raw bytes incurs excessive computational expense owing to increased sequence lengths~\citep{xue2022byt5}. Earlier approaches address this issue by adopting more resource-efficient self-attention techniques~\citep{el2020characterbert, clark2022canine} or by leveraging architectures that do not rely on attention mechanisms~\citep{wang2024mambabyte}~(\S\ref{section:related_work}). Nevertheless, these strategies are chiefly advantageous for training \textit{smaller models}. When scaled up, the dominant factor in the computational burden of a Transformer is the sizable feed-forward network layers that must be executed for each byte, while the attention component becomes less significant in terms of overall cost.

In order to optimize compute allocation, we introduce a flexible, adaptive approach that organizes bytes into \textit{patches}~(\S\ref{section:patching}) and present a novel model architecture capable of integrating both byte-level and patch-level information. In contrast to conventional tokenization, {\textsc BLT} does not impose a predefined set of patch vocabulary. Instead, any combination of bytes can be transformed into latent patch embeddings using compact, trainable encoder and decoder components. Our findings demonstrate that this method allows for \textit{greater} computational efficiency compared to traditional tokenization-based approaches.

LLMs that rely on tokenization assign an equal computational budget to every token, regardless of the underlying complexity. This compromises computational efficiency in favor of uniform processing, as token boundaries—determined by compression strategies—may not reliably reflect the intricacy involved in prediction tasks. Our architectural approach is rooted in the principle that computation should be adaptively targeted to regions that require it most. For instance, predicting the final characters in most words generally constitutes a simpler, lower-entropy task than generating the initial token of a new sentence, and therefore does not warrant the involvement of a large transformer. This philosophy is embodied in {\textsc BLT}'s design~(\S\ref{section:architecture}), which consists of three transformer modules: two lightweight, byte-level \textit{local models} complemented by a larger, global \textit{latent transformer}~(\autoref{fig:arch_high_level}). 

To facilitate dynamic allocation of computational resources and determine the appropriate grouping of bytes into patches, {\textsc BLT} partitions the data according to the entropy measured in next-byte prediction, thus yielding context-aware clusters of bytes that maintain similar information density.

We conduct the first flop-controlled scaling analysis of byte-level models reaching up to 8B parameters and trained on 4T bytes, demonstrating that it is feasible to train large-scale models from raw bytes directly, foregoing the use of fixed-vocabulary tokenization. In summary, {\textsc BLT} achieves parity with Llama 3 in terms of flop-controlled training performance\footnote{The computational budget for models is assessed via the tally of Floating Point OPerations (flops) required.}, yet requires as much as 50\% fewer flops during inference~(\S\ref{section:scaling}). Our results also reveal that handling data at the byte level leads to marked gains when modeling the infrequent cases in the dataset. The {\textsc BLT} models exhibit superior robustness to noisy input relative to models employing tokenizers, and they possess stronger abilities in character-level comprehension as evidenced by improved outcomes on tests of orthographic processing, phonological understanding, and machine translation for low-resource languages~(\S\ref{sec:robustness}). Moreover, with the {\textsc BLT} framework, both model size and patch size can be expanded concurrently without exceeding the preset inference flop constraint. Utilizing longer patch sizes generally reduces computation demands, enabling those savings to be reinvested in increasing the global latent transformer’s size, since it operates less frequently. Through experiments that maintain a fixed inference flop allowance~(\autoref{fig:fixed_inference_scaling}), we detect considerably improved scaling characteristics compared to approaches that rely on tokenization.

To conclude, the primary contributions of this work are as follows: 1) We present {\textsc BLT}, a byte latent LLM framework designed to dynamically manage computation, thereby enhancing flop efficiency; 2) We provide evidence that {\textsc BLT} achieves parity with Llama 3, in terms of flop-controlled training, up to the 8B parameter level, while also offering the ability to exchange slight declines in evaluation scores for improvements in flop efficiency reaching as high as 50\%; 3) {\textsc BLT} enables an alternative scaling paradigm for LLMs, making it possible to increase model size while keeping inference cost constant; 4) We establish that {\textsc BLT} exhibits superior resilience to noisy inputs and demonstrates sensitivity to sub-word elements within the input data, which are typically overlooked by token-based LLMs. The training and inference implementation for {\textsc BLT} has been made publicly accessible at~\url{https://github.com/facebookresearch/blt}.

\section{Patching: From Individual Bytes to Groups of Bytes}
\label{section:patching}

Dividing byte sequences into \textit{patches} enables {\textsc BLT} to flexibly assign computational resources in response to varying contextual requirements. Figure~\ref{fig:patching_types} illustrates multiple strategies for organizing bytes into patches. More specifically, a patching function $f_p$ partitions a byte sequence $\pmb{x}=\{x_i,|i=1,\ldots n\}$ of length $n$ into a collection of $m < n$ patches, denoted $\pmb{p}=\{p_j|j=1,\ldots,m\}$, by mapping each element $x_i$ to \{0,1\}, where a value of 1 signifies the beginning of a new patch. For both token-oriented and patch-oriented architectures, the dominant factor in computational overhead arises from the number of steps the main Transformer executes, which in {\textsc BLT} is dictated by the number of patches created by a particular patching function. As a result, the average number of bytes in a patch—or \textit{patch size}—serves as the principal metric governing the compute cost associated with data processing during both model training and inference when a given patching approach is employed~(\S\ref{section:flops}).
We proceed by describing three specific patching functions: segmenting with a constant byte count per patch~(\S\ref{section:static-patch}), using whitespace as a delimiter for patching~(\S\ref{section:white-patch}), and adaptively forming patches based on local entropy estimates from a compact byte-level language model~(\S\ref{section:dyn-patch}). In conclusion, we provide an overview of incremental patching, and clarify the distinction between patching and tokenization~(\S\ref{section:bpe-patch}).

\subsection{Strided Patching Every K Bytes}
\label{section:static-patch}

One of the simplest approaches to grouping bytes is to segment them into fixed-size patches of length $k$, following the method employed in MegaByte~\citep{yu2023megabyte}. This approach is advantageous in that a constant stride is straightforward to implement for both model training and inference, and it facilitates easy adjustment of the average patch length, granting explicit control over computational cost in terms of flops. Nonetheless, this patching scheme introduces notable limitations. Primarily, computation is not adaptively distributed according to the local content: transformer steps may be wasted on regions of low complexity, such as runs of whitespace in source code, while regions with a high density of semantic content—like mathematical notation—might not receive sufficient computational attention. Additionally, fixed-size patching can produce inconsistencies by segmenting identical byte sequences, such as repeated words, in varying and non-contextual ways.

\subsection{Space Patching}
\label{section:white-patch}

\citet{slagle2024spacebyte} introduces an efficient modification to strided patching that forms new patches at each occurrence of space-like bytes\footnote{
    Space-like bytes refer to any byte that is neither a Latin character, digit, nor a UTF-8 continuation byte. Additionally, every patch must include at least one non space-like byte. }—positions that often correspond to linguistic boundaries in numerous languages. In the approach termed Space patching, every word is given a dedicated latent transformer step (involving increased flops), ensuring a consistent patching scheme for words across different sequences and directing computational resources towards more challenging predictions, which typically appear after spaces. As an illustrative case, the task of generating the initial byte in the answer to “Who composed the Magic Flute?\uline{\hspace{.6em}}” is considerably more complex than generating subsequent bytes after “M,” as the choice for the first character dramatically narrows the set of plausible completions, making the remainder of “Mozart” easier to infer. Nonetheless, the space patching technique exhibits limitations—it does not generalize smoothly to every language or domain and lacks the flexibility to adapt patch sizes. In the following, we present an alternative patching strategy informed by the observation that word-initial bytes pose the greatest prediction challenge, and that inherently offers a way to adjust patch size as needed.

\subsection{Entropy Patching: Using Next-Byte Entropies from a Small Byte LM}
\label{section:dyn-patch}

Instead of depending on heuristic rules like whitespace for boundary detection, we adopt a data-centric methodology to pinpoint next-byte predictions with significant uncertainty. We propose \textit{entropy patching}, a technique that leverages entropy estimates to determine patch boundaries.

To achieve this, we train a compact byte-level auto-regressive language model using the training set for {\textsc BLT}. This model is then utilized to calculate the entropy of the next byte under the language model probability distribution $p_{e}$ across the byte vocabulary $\mathcal{V}$:

\begin{equation}
H(x_i) = \sum_{v \in \mathcal{V}} p_{e}(x_i=v|\pmb{x}_{<i}) \log p_{e}(x_i=v|\pmb{x}_{<i})
\end{equation}

We explore two approaches for determining patch boundaries using entropies $H(x_i)$. The initial method selects locations where the entropy surpasses a fixed global threshold, as depicted in~\autoref{fig:patching}. In contrast, the second technique highlights positions exhibiting a sharp increase in entropy compared to the immediately preceding point. This latter strategy can also be viewed as pinpointing locations where the roughly monotonic decrease in entropy within a patch is disrupted.

\begin{align*}
    \text{Global Constraint}\hspace{1cm}H(x_t)& > \theta_g\\
    \text{Approx. Monotonic Constraint}\hspace{1cm}H(x_t)& - H(x_{t-1}) > \theta_r
\end{align*}

Patch boundaries are determined through a minimal preprocessing procedure that takes place as part of the dataloading process. This contrasts with the approach of~\citet{nawrot-etal-2023-efficient}, where a classifier is employed to infer patch boundaries based on entropy measurements. In our study~(\S\ref{section:experiments}), we evaluate and contrast these two strategies for discriminating between bytes of low and high entropy.

\subsection{The Byte-Pair Encoding (BPE) Tokenizer and Incremental Patching}
\label{section:incremental-patching}
\label{section:bpe-patch}

A wide range of contemporary llms, such as our baseline Llama 3, employ subword tokenizers like bpe~\citep{gage1994new, sennrich-etal-2016-neural}. In our context, ``tokens'' denote byte-groupings selected from a \textit{finite} vocabulary established before training begins, in contrast to ``patches,'' which represent variable-length groupings generated on the fly and lack a predefined vocabulary. One key distinction between tokens and patches is that when working with tokens, the model does not have direct exposure to the raw byte-level information.

An important advancement introduced by {\textsc BLT} relative to tokenization-based architectures is its reconfiguration of the balance between vocabulary size and computational requirements. In conventional large language models (LLMs), expanding the vocabulary typically results in tokens that encapsulate more information on average, thereby reducing the total number of inference steps needed; however, this simultaneously enlarges the dimensionality of the model's final output projection. This inherent compromise restricts how much token-based models can meaningfully alter token lengths or inference efficiency. For instance, Llama 3 increases the mean token size from 3.7 to 4.4 bytes, yet this enhancement comes at the expense of quadrupling its embedding table in comparison to Llama 2.

During generation, {\textsc BLT} must determine at each point in the byte sequence whether it has reached a patch boundary, as this influences whether additional computation is performed using the Latent Transformer. Importantly, this determination must be made without knowledge of subsequent, as-yet-unproduced portions of the sequence. Therefore, the method of segmenting the byte sequence into patches must operate solely with access to the already generated prefix and cannot depend on future bytes. More precisely, a patching method $f_p$ exhibits incremental patching if:
$$f_p(\pmb{x}_{<i}) = f_p(\pmb{x})_{<i}$$
bpe does not fulfill the requirements of incremental patching, since a given prefix may be tokenized differently contingent on its subsequent continuation. Hence, bpe violates the above property.\footnote{Introducing a special delimiter token to explicitly mark patch boundaries allows bpe to conform to incremental patching, but this comes at the cost of longer byte sequences.}

\section{{\textsc BLT} Architecture}
\label{section:architecture}
The {\textsc BLT} system consists of a large, global autoregressive language model that operates at the patch level, supplemented by two lightweight local models responsible for mapping byte sequences to patches and reconstructing bytes from patch representations~(\autoref{fig:arch_high_level}).

\subsection{Latent Global Transformer Model}

\textit{The Latent Global Transformer} refers to an autoregressive transformer architecture, denoted by $\mathcal{G}$ and comprising $l_{\mathcal{G}}$ layers, that transforms a sequence of latent input patch embeddings $p_j$ into a sequence of corresponding output patch embeddings $o_j$. In this work, indices $j$ and $i$ are consistently used to indicate patches and bytes, respectively. The global transformer incorporates a block-causal attention mask~\citep{dubey2024llama}, limiting each patch's attention span to itself and preceding patches within a document. The computational demands of this global model dominate both pre-training and inference stages. Consequently, the frequency with which it is employed can be modulated to allocate computational resources adaptively based on the complexity of specific input and output segments.

\subsection{Local Encoder}
\label{section:local}

\textit{The Local Encoder Model}, represented by $\mathcal{E}$, comprises a compact transformer architecture utilizing $l_{\mathcal{E}}$ layers, where $l_{\mathcal{E}} << l_{\mathcal{G}}$. Its primary function is to rapidly convert sequences of input bytes $b_i$ into rich patch-level representations $p_j$. Distinguishing this design from the canonical transformer, a cross-attention mechanism is inserted subsequent to each transformer layer, aggregating byte-level information into patch-level features as illustrated in~(\autoref{fig:crossattn}).

The model first maps each byte $b_i$ to an embedding $x_i$ using an embedding matrix of size $\mathbb{R}^{256 \times h_{\mathcal{E}}}$. These embeddings can optionally be enriched with hash-embeddings for supplementary contextual signals~(\S\ref{sec:hash-emb}). Through multiple layers consisting of alternating transformer blocks and cross-attention modules~(\S\ref{sec:enc-cross-attn}), the architecture transforms byte representations into patch representations, $p_i$, which serve as input to the global transformer $\mathcal{G}$.

Within the transformer blocks, a \textit{local block causal} attention scheme is utilized: each token is restricted to attend to a sliding window of $w_{\mathcal{E}}$ preceding bytes. This attention window is permitted to cross dynamically-defined patch boundaries, yet is not allowed to span across separate document segments. The forthcoming sections detail the embedding process and elucidate the implementation of the cross-attention component.

\subsubsection{Encoder Hash n-gram Embeddings}
\label{sec:ngram-emb}
\label{sec:hash-emb}
An essential aspect of producing robust and informative representations at each position $i$ is integrating knowledge of the prior bytes. In {\textsc BLT}, this is accomplished by representing the byte $b_i$ on its own, as well as within the context of an n-gram of bytes. Specifically, at each position $i$, we begin by generating byte-grams

\begin{align}
    g_{i,n}=\{b_{i-n + 1},\ldots, b_i\}
\end{align}

for every byte position $i$ and for $n$ values ranging from three to eight.\footnote{
Byte-grams with length $n$ or greater are excluded in cases where $i < n$. }

Next, we propose the use of hash $n$-gram embeddings: each byte $n$-gram is assigned an index within a fixed-size embedding table $E_{n}^{hash}$ through the application of a hash function, independently for each $n \in\{3,4,5,6,7,8\}$~\citep{bai2010learning}.
The corresponding $n$-gram embedding is combined additively with the original byte embedding, then normalized before being supplied as input to the local encoder. The enriched embedding is computed as

\begin{align}
e_i &= x_i + \sum_{n = 3,...,8}E_{n}^{hash}(\text{Hash}(g_{i,n})) \\
\text{where, Hash}(g_{i,n}) &= \text{RollPolyHash}(g_{i,n}) \% |E_{n}^{hash}|
\end{align}

We divide $e_i$ by the total count of $n$-gram sizes incremented by one, applying RollPolyHash as described in Appendix~\ref{appendix:rollpolyhash}. Section~\ref{section:ablations} presents an ablation study assessing the impact of $n$-gram hash embeddings under varying $n$ values and embedding table sizes on trends governed by flop-controlled scaling laws. Alongside hash $n$-gram embeddings, our investigations included frequency-based $n$-gram embeddings; further information regarding this line of inquiry can be found in Appendix~\ref{appendix:ngrams}.

\subsubsection{Encoder Multi-Headed Cross-Attention}
\label{sec:enc-cross-attn}
Our approach is largely inspired by the cross-attention mechanism used in the Perceiver architecture~\citep{jaegle2021perceiver}, with the key distinction that, rather than relying on a fixed array of latent variables, our latent representations are derived from variable patch embeddings~(\autoref{fig:crossattn}), and the attention operation is restricted to only those bytes within each individual patch. The mechanism constructs a query vector for every patch $p_j$, initialized by aggregating (pooling) the byte representations associated with that particular patch $p_j$, then applying a linear transformation, $\mathcal{E}_{C} \in \mathbb{R}^{h_{\mathcal{E}} \times (h_{\mathcal{E}} \times U_{\mathcal{E}})}$, where $U_{\mathcal{E}}$ represents the total number of encoder cross-attention heads. Explicitly, for patch $p_j$ with its associated byte sequence $f_{\text{bytes}}(p_j)$, the computation proceeds as follows

\begin{align}
P_{0,j} &= \mathcal{E}_{C}(f_\text{bytes}((p_j)), f~ \text{is a pooling function} \\
P_l &= P_{l-1} + W_o\left(\text{softmax}\left(\frac{QK^T}{\sqrt{d_k}}\right)V\right) \\
\text{where } Q_j &= W_q(P_{l-1,j}), K_i = W_k(h_{l-1,i}), V_i = W_v(h_{l-1,i}) \\
h_l &= \text{Encoder-Transformer-Layer}_l(h_{l-1})
\end{align}

Here, $P \in \mathbb{R}^{n_p \times h_{\mathcal{G}}}$ denotes the set of $n_p$ patch representations that serve as input to the global model. The initial representation of each patch, $p_j$, is obtained by pooling the byte embeddings $e_i$ associated with it. The matrices $W_q$, $W_k$, $W_v$, and $W_o$ perform linear projections for queries, keys, values, and outputs, respectively, where keys and values are computed from the byte-level representations $h_i$ of the previous layer (using $e_i$ in the first layer). We adopt a patch-specific masking approach so that each query $Q_j$ only considers keys and values corresponding to the bytes contained within the same patch $j$. Since our architecture employs multi-head attention on $Q, K$, and $V$ and because patch-level embeddings generally have a higher dimensionality ($h_{\mathcal{G}}$) compared to $h_{\mathcal{E}}$, during cross-attention we keep $P_l$ as multiple attention heads with dimension $h_{\mathcal{E}}$, subsequently concatenating them to form $h_{\mathcal{G}}$-dimensional vectors. Moreover, we apply a pre-LayerNorm operation to queries, keys, and values, omitting any use of positional embeddings in the cross-attention block. The output of the cross-attention layer is combined with its input via a residual connection.

\subsection{Local Decoder} 

The local decoder $\mathcal{D}$, like its encoder counterpart, utilizes a streamlined transformer-based architecture characterized by a relatively small number of layers, $l_{\mathcal{D}} << l_{\mathcal{G}}$. Its primary function is to convert the global patch representations $o_j$ back into sequences of raw bytes, denoted as $y_i$. During the decoding process, it generates each byte based on the context provided by the bytes decoded so far, leveraging the hidden states output by the local encoder corresponding to the input byte sequence. The architecture interleaves $l_{\mathcal{D}}$ layers composed alternately of cross-attention and transformer operations. Initially, the cross-attention mechanism transforms patch representations into representations suitable for bytes, after which the transformer layer of the decoder refines these byte-level sequences.

\subsubsection{Decoder Multi-headed Cross-Attention} 

Within the decoder’s cross-attention mechanism, the positions of queries and key/values are swapped: here, byte-representations serve as queries, whereas patch representations take on the role of key/values. To initialize the cross-attention, the byte-representations are set using the byte embeddings produced by the final encoder layer, denoted as $h_{l_{\mathcal{E}}}$. For each following decoder layer $l$, the updated byte-representations $d_{l,i}$ are determined by:

\begin{align}
D_0 &= h_{l_{\mathcal{E}}} \\
B_l &= D_{l-1} + W_o\left(\text{softmax}\left(\frac{QK^T}{\sqrt{d_k}}\right)V\right), \\
  &\text{where } Q_i = W_q(d_{l-1,i}), K_i = W_k(\mathcal{D}_C(o_j)), V_i = W_v(\mathcal{D}_C(o_j)) \\
  D_l &= \text{Decoder-Transformer-layer}_l(B_l)
\end{align}

Here, $W_k$ and $W_v$ denote the key and value projection matrices, which are used to perform a linear transformation followed by the split operation $\mathcal{D}_C$ on the final patch representations $o_j$ produced by the global model. The matrix $W_q$ is responsible for projecting queries, and it acts on the byte-level representations $d_{l-1}$ from the preceding decoder transformer layer (or on $h_{l_{\mathcal{E}}}$ when initializing the first layer). $W_o$ serves as the output projection matrix. As a result, $B$ has the shape $\mathbb{R}^{h_{\mathcal{D}} \times n_b}$, where $n_b$ represents the total number of output bytes. The subsequent decoder representations $D_l$ are generated by applying a decoder transformer layer to the output $B$ from the cross-attention module. Consistent with the approach used in local encoder cross-attention, we utilize multiple attention heads, implement pre LayerNorms, exclude positional embeddings, and include a residual connection that wraps around the cross-attention block.

\section{Experimental Setup}
\label{section:experiments}
We systematically construct controlled experiments to compare the performance of {\textsc BLT} against models that rely on tokenization, ensuring that {\textsc BLT} does not receive any benefits that could arise from the ability to process longer sequence contexts.

\subsection{Pre-training Datasets} In this work, all model variants are initially trained on two main datasets: (1) The Llama 2 dataset~\citep{touvron2023llama}, consisting of 2 trillion tokens drawn from a wide range of publicly accessible resources, meticulously cleaned and filtered to ensure higher quality; and (2) {\textsc BLT}-1T: a newly assembled corpus containing 1 trillion tokens, curated from assorted open-access sources, which also incorporates a portion of the pre-training data publicly provided by Datacomp-LM~\citep{li2024datacomp}. The first dataset is leveraged for scaling law analyses aimed at identifying the optimal token budget as prescribed by~\cite{dubey2024llama}, informing key architectural parameters for {\textsc BLT}. The second dataset serves as the foundation for full-scale pre-training, enabling performance comparison with Llama 3 across various downstream evaluations. Importantly, neither collection contains any content sourced from Meta’s proprietary products or platforms. Additionally, for baseline comparisons among models using tokenizers, we adopt the Llama 3 tokenizer, which features a vocabulary of 128K tokens and yielded better baseline outcomes in our trials relative to the Llama 2 tokenizer.

\subsection{Entropy Model} For our experiments, the entropy model is implemented as a byte-level language model, which is trained on the identical data distribution as the comprehensive {\textsc BLT} model. By default, the architecture consists of a transformer containing 100M parameters, organized into 14 layers, with a hidden size of 512 and a sliding window attention mechanism covering 512 bytes. All other hyperparameter settings mirror those used in our local and global transformer configurations. We also performed a series of investigations into various model capacities, receptive field sizes, and alternative architectures, as elaborated in \autoref{section:ablations}. Notably, if the model's receptive field is sufficiently restricted, the resulting entropy model can be stored as an efficient lookup table.

\subsection{Entropy Threshold and Equalizing Context Length} For entropy-based patching models, we determine a patching threshold that yields a target average \textit{patch size} across the pretraining data mix. In the context of {\textsc BLT}, in contrast to tokenization, the \textit{patch size} is flexibly set, which substantially affects the resulting context length available to the model. To standardize the average context length and prevent models with larger patch sizes from gaining an undue advantage, we constrain the total number of bytes per batch to be the same on average. As a result, for configurations with larger patch sizes, the sequence length is decreased. Specifically, when training on Llama 2 data, we adopt an 8k byte context, whereas with the {\textsc BLT}-1T dataset, the context is expanded to 16k bytes on average. In both cases, the average batch size is kept fixed at 16M bytes.

Although the average batch size remains unchanged, dynamic patching approaches result in varying proportions of bytes per patch across batches. To maximize efficiency, our implementation of {\textsc BLT} training organizes batches such that patches are densely packed, thereby eliminating the need for padding in the computationally intensive latent transformer. This packing guarantees a uniform patch count across all batches. Throughout training, we pad or truncate byte sequences to lengths of 12k for Llama 2 and 24k for the {\textsc BLT}-1T datasets, as necessary, in order to prevent unexpected memory surges caused by exceptionally large patches.

\subsection{Entropy Model Context}
\label{section:entropy-context}
Our experiments indicate that employing entropy patching leads to increasingly larger patches, particularly in structured domains such as multiple choice formats, which frequently contain repetitive segments (refer to \autoref{fig:mmlu_patching} for an MMLU illustration). These discrepancies arise due to reduced entropy associated with the recurrent portions present in the entropy model context. Therefore, in the primary experiment involving {\textsc BLT}-Entropy with a patch size of 4.5, we reinitialize the entropy context by inserting line breaks and implement an approximate monotonicity constraint, as this approach is less susceptible to "entropy drift" resulting from fluctuations in context length. Notably, this adjustment solely impacts the manner in which entropies are computed; the method for determining the entropy threshold value remains consistent.

\subsection{FLOPs Estimation} 
\label{section:flops}

Our approach to calculating transformer flops is primarily based on the methodology presented in Chinchilla~\citep{hoffmann2022training}, which includes flops for the feed-forward modules, the qkvo projections within self-attention, as well as for the attention calculation and output projection. One important deviation is our assumption that the input embedding layer functions as an efficient lookup table, rather than relying on dense matrix multiplication, thereby reducing its flop count to zero. In line with earlier studies, we approximate the computational cost of the backward pass as being twice that of the forward pass.

To determine the number of flops \textit{per byte} for {\textsc BLT} models, we aggregate the flop counts from several sources: the local encoder transformer, the global latent transformer, and the local decoder transformer, as well as the contributions from the cross-attention modules in both the encoder and decoder:

\begin{align}
    \text{FL}_{\text{{\textsc BLT}}} &= \text{Transf. FL}(h_{\mathcal{G}}, l_{\mathcal{G}}, m=n_{ctx}/n_p, V=0) / n_p\\
    &\quad +\text{Transf. FL}(h_{\mathcal{E}}, l_{\mathcal{E}}, m=w_{\mathcal{E}}, V=0) \\
    &\quad + \text{Transf. FL}(h_{\mathcal{D}}, l_{\mathcal{D}}, m=w_{\mathcal{D}}, V=256) \\
    &\quad + \text{Cross Attn. FL}(h_{\mathcal{E}}, l_{\mathcal{E}}, m=n_p, r=n_p/k) \times k/n_p \\
    &\quad + \text{Cross Attn. FL}(h_{\mathcal{D}}, l_{\mathcal{D}}, m=k, r=k/n_p)
\end{align}

In this context, $n_{ctx}$ represents the length of the sequence measured in bytes, while $n_p$ denotes the patch size. The parameter $r$ specifies the proportion of queries relative to keys and values. The value $k$ is defined as the ratio between the patch dimension and the byte dimension, indicating how many segments of the local model are concatenated to produce the full model representation (with $k=2$ as illustrated in \autoref{fig:crossattn}). The term $V$ stands for the size of the output vocabulary and is relevant exclusively to the local decoder. The attention mechanism adopts a specific context length, $m$, which varies depending on whether the module processes a byte-level or a patch-level sequence. We adjust the calculation of attention flops for each distinct component. Precise formulas detailing the flops computations for both Transformer and Cross-Attention modules are included in Appendix~\ref{section:flops-appendix}.

\subsection{Bits-Per-Byte Estimation} The metric of perplexity is inherently linked to the choice of tokenizer, as it reflects model uncertainty on a per-token basis. To facilitate a fair comparison between byte-level and token-level approaches, we adopt Bits-Per-Byte (BPB)—an alternative that is agnostic to the tokenizer and has been utilized in earlier research~\citep{xue2022byt5, yu2023megabyte, wang2024mambabyte}. BPB is computed as follows:

\begin{align}
    \text{BPB}(x)&= \frac{\mathcal{L}_{CE}(\pmb{x})}{\ln(2)\cdot n_{\text{bytes}}}
\end{align}

Here, the cumulative cross-entropy loss over the input $\pmb{x}$ is scaled by both the total byte count in $\pmb{x}$ and a normalization factor, yielding a standardized uncertainty measure per byte.

\subsection{Transformer Architecture Hyperparameters} 

In all transformer blocks within {\textsc BLT}, encompassing both local and global components, our approach is primarily modeled after the design of Llama~3~\citep{dubey2024llama}. Specifically, feed-forward layers employ the SwiGLU activation~\citep{swiglu}, and rotary positional embeddings (RoPE)~\citep{RoPe}—with $\theta=500000$ as suggested by~\citep{xiong2024effective}—are incorporated exclusively into self-attention modules. For normalization, we utilize RMSNorm~\citep{RMSNorm} throughout the network. Flash attention~\citep{dao2022flashattention} is applied across all self-attention layers configured with fixed-standard attention masks, such as \textit{block causal} or \textit{fixed-window block causal} schemes, and we set the attention mask’s window size to 512 in fixed-width scenarios. Because cross-attention layers require masks that adapt dynamically based on patches, we leverage Flex Attention\footnote{\url{https://pytorch.org/blog/flexattention}}, which provides optimized fused kernels, thereby markedly accelerating training.

\subsection{{\textsc BLT}-Specific Hyperparameters} To assess the performance of {\textsc BLT} models, we perform experiments in two primary directions: examining scaling laws and evaluating downstream task performance, using models of varying sizes—specifically 400M, 1B, 2B, 4B, and 8B parameters. Details regarding the architectural hyperparameters for each model are compiled in Appendix~Table~\ref{tab:arch}.

The initialization of queries for the first cross-attention layer in the local encoder employs max-pooling. Across all {\textsc BLT} models, we utilize $500,000$ hashes with a single hash function, adopting n-gram lengths spanning from 3 to 8. The learning rate is consistently set at $4e-4$ for each model. This uniform learning rate selection for both token and {\textsc BLT} models is informed by a hyperparameter search within the $1e-3$ to $1e-4$ interval at the 400M and 1B parameter settings, which revealed the shared optimal value. When analyzing scaling on the Llama-2 dataset, batch sizes are determined following the recommendations from~\cite{dubey2024llama}, or based on byte-wise equivalents.

For the optimization process, we adopt the AdamW optimizer~\citep{adamw}, configuring $\beta_1$ to 0.9 and $\beta_2$ to 0.95, with $\epsilon=10^{-8}$. The training schedule includes a linear warm-up phase over 2000 steps, after which the learning rate decays to zero following a cosine schedule. Weight decay is set at 0.1, and global gradient norms are clipped at a maximum value of 1.0.

\section{Scaling Trends}
\label{section:scaling}

We offer a comprehensive analysis of scaling behaviors observed in byte-level models, with a focus on informing future expansion of {\textsc BLT} models. Our investigation seeks to overcome shortcomings found in earlier studies of byte-level modeling through several key contributions: (a) conducting direct comparisons under compute-optimal training paradigms, (b) developing 8B parameter models trained with substantial datasets—reaching up to 1T tokens or 4T bytes—and assessing their performance on a range of downstream benchmarks, and (c) systematically examining scaling effects under conditions where inference costs are explicitly managed. Subsequently, we will explore distinctive benefits that arise from the modeling of byte-level sequences.

\subsection{Parameter Matched Compute Optimal Scaling Trends}
\label{section:parameter-matched-scaling}
We utilize the Llama 2 dataset to train a range of \textit{compute-optimal} bpe and {\textsc BLT} models, each at four distinct scales from 1B up to 8B parameters. To assess scaling behavior, we chart the relationship between training flops and language modeling accuracy on a representative subset of the training data mixture. The bpe models are optimized using the parameter-to-dataset size ratio recommended in Llama~3~\citep{dubey2024llama}, which is constructed to optimize performance within a given computational training constraint~\citep{hoffmann2022training}. This approach yields a strong baseline for comparison. For every bpe model configuration, an identically sized and architecturally matched {\textsc BLT} variant is also trained on the same data, employing a Latent Transformer with the same underlying architecture as its bpe counterpart.

As shown in~\autoref{fig:scaling-compute-opt} (right), {\textsc BLT} models consistently equal or surpass the performance of their bpe equivalents, a pattern that persists across increasing model sizes and flops. According to our knowledge, {\textsc BLT} represents the first byte-level Transformer architecture to demonstrate comparable scaling behaviors to BPE-based models under compute-optimal conditions. These results thus support our hypothesis that the optimal parameter-to-compute ratio established for bpe is similarly applicable to {\textsc BLT}, or at the very least, does not differ substantially.

Both enhancements to the architecture and the incorporation of dynamic patching are essential for aligning with the scaling behavior of bpe. In ~\autoref{fig:scaling-compute-opt} (left), we present a comparison between space-patching models and Llama 3. To approximate SpaceByte \citep{slagle2024spacebyte}, we employ {\textsc BLT} space-patching without utilizing n-gram embeddings or cross-attention mechanisms. While SpaceByte achieves better results than Megabyte, there is still a significant performance gap compared to Llama 3. The right side of \autoref{fig:scaling-compute-opt} highlights the collective benefits brought about by architectural innovations alongside dynamic patching. At large scales, {\textsc BLT} models demonstrate competitive performance with leading tokenizer-based models such as Llama 3.

Furthermore, our findings indicate that, among tokenizer-based approaches, model performance is influenced by tokenizer selection. Specifically, models utilizing the Llama-3 tokenizer exhibit superior outcomes relative to those employing the Llama-2 tokenizer when trained on identical data.

Lastly, our {\textsc BLT} architecture exhibits performance characteristics that interpolate between those of Llama 2 and Llama 3 when much larger patch sizes are employed. While the average token produced by the bpe tokenizers in Llama 2 and 3 measures 3.7 and 4.4 bytes, respectively, {\textsc BLT} demonstrates analogous scaling behavior with average patch sizes of 6 or even 8 bytes. Since inference flop is inversely related to average patch size, setting the patch size to 8 bytes can nearly halve the computational requirements for inference. Moreover, as both model and dataset sizes increase, models utilizing larger patch sizes tend to yield superior outcomes. For instance, although {\textsc BLT} configured with an 8-byte patch size initially underperforms relative to the bpe-tokenized Llama 2 at the 1B parameter scale, it surpasses bpe-based models at the 7B scale. This observation indicates that larger patch sizes may confer increasing benefits at greater scales, and suggests the potential viability of even greater patch sizes as model capacity and available training compute expand.

\subsection{Beyond Compute Optimal Task Evaluations}
\label{section:evals}
To further explore model scaling behaviors, we extend training of an 8B {\textsc BLT} model past the compute-optimal ratio using the {\textsc BLT}-1T dataset, which is both larger and higher in quality. We then evaluate its performance using a comprehensive set of standard classification and generation tasks. The chosen benchmarks encompass areas such as common sense reasoning, world knowledge, and code generation:
\paragraph{Classification tasks} We assess performance on ARC-Easy (0-shot)~\citep{Eval_ARC}, Arc-Challenge (0-shot)~\citep{Eval_ARC}, HellaSwag (0-shot)~\citep{Eval_hellaswag}, PIQA (0-shot)~\citep{Eval_piqa}, and MMLU (5-shot)~\citep{hendrycks2020measuring}. For these tasks, a prompt-based scoring approach is utilized by computing likelihoods over the possible answer choices; final results are presented as mean accuracies.
\paragraph{Coding related generation tasks:} Pass@1 results are provided for both MBPP (3-shot)~\citep{Eval_MBPP} and HumanEval (0-shot)~\citep{Eval_HumanEval} to measure large language models’ Python code generation capabilities.

In \autoref{tab:evals}, we present a comparison among three models, each trained using the {\textsc BLT}-1T dataset: a model based on the bpe Llama 3 tokenizer,\footnote{We opt for the Llama 3 tokenizer with its 128k vocabulary as it demonstrates superior performance relative to Llama 2's 32k vocabulary.} along with two distinct variants of the {\textsc BLT} architecture. Specifically, one variant adopts a space-patching approach ({\textsc BLT}-Space), while the other leverages an entropy-driven patching mechanism ({\textsc BLT}-Entropy), incorporating an approximate monotonicity constraint and resetting the entropy model’s context at each newline event, as elaborated in~\autoref{section:entropy-context}. The training of all three models is carried out using a matched flop budget. Notably, for the {\textsc BLT}-Entropy model, we modify the entropy threshold during inference from 0.6 to 0.1; this tuning results in enhanced task performance, albeit with a trade-off of increased inference steps.

The {\textsc BLT}-Entropy model achieves superior results compared to the Llama 3 model on 4 out of the 7 evaluated tasks, despite both models being trained on an equal number of bytes. This enhanced performance can be attributed to two main factors: (1) more efficient utilization of training computation enabled by dynamic patching, and (2) the explicit modeling of information at the byte level rather than relying on token-based representations.

Conversely, {\textsc BLT}-Space lags behind the Llama 3 tokenizer in performance across nearly all tasks, except for one. However, thanks to its higher average patch size of 6 bytes, it is able to significantly decrease inference flops. By contrast, models based on bpe and entropy-patching strategies exhibit similar average patch sizes of about 4.5 bytes on the mixed training dataset. Given identical training resources, the model employing a larger patch size can process 30\% more data than the other two approaches, potentially moving {\textsc BLT} farther from the compute-optimal threshold.

\subsection{Patches Scale Better Than Tokens} In {\textsc BLT} models, it is possible to expand both the model's capacity and the patch dimensions concurrently, all while holding the training and inference flop counts as well as the volume of training data steady. Unlike traditional token-based models that rely on a fixed vocabulary, patch-based approaches offer the distinctive capability to arbitrarily enlarge patch size, circumventing the inherent efficiency limitations described in Section~\ref{section:bpe-patch}. Increasing the patch length directly reduces computational demands, enabling more resources to be devoted to scaling up the global latent transformer, as its execution frequency decreases.

A fixed inference scaling study is performed to evaluate whether increasing model size and reducing the number of inference steps on larger patches leads to superior performance compared to smaller models that process more steps. Using initial model configurations of 400 million and 3.6 billion parameters with the Llama 2 tokenizer, we identify flop-equivalent architectures employing the Llama 3 tokenizer and {\textsc BLT}-Entropy models, which utilize average patch sizes of 6 and 8 bytes on the training datamix (refer to~\autoref{tab:fixed-inf-params} for detailed model specifications). For models operating with a patch size of 8, the architecture includes 3 encoder layers rather than a single layer. All models are trained under a variety of predefined training flop budgets.

\autoref{fig:fixed_inference_scaling} illustrates that {\textsc BLT} models exhibit superior scaling behavior compared to tokenization-based counterparts across both categories of inference computational cost. Specifically, BPE models demonstrate an initial advantage under limited training resources, yet are overtaken by {\textsc BLT} models shortly after exceeding the compute-optimal threshold. From a practical standpoint, allocating greater resources to one-time pretraining can yield models with improved performance under a constant inference budget. This pattern is exemplified by 8B parameter models such as Llama 3.1, which have been pretrained on datasets nearly two orders of magnitude larger than what would be optimal purely from a compute efficiency perspective for their model size.

The point at which {\textsc BLT} surpasses token-based approaches shifts marginally nearer to the compute-optimal threshold as model size increases, moving from approximately three times to 2.5 times the optimal compute budget in larger flop regimes. Likewise, the model employing patch size 8 demonstrates a more pronounced scaling behavior in these higher flop categories, overtaking alternatives more rapidly. As highlighted in Section~\ref{section:parameter-matched-scaling}, employing larger patch sizes leads to performance that increasingly approximates BPE models at greater model scales. We partially attribute this to the reduced proportion of overall flops consumed by the byte-level Encoder and Decoder components, as their scaling rate appears slower than that of the Latent Transformer. When increasing the model’s total parameter count by a factor of 20 (from 400M to 8B), the local parameters specific to {\textsc BLT} are only about doubled. This observation matters because adjustments to patch size influence only the flops allocated to the patch Latent Transformer, not those of the byte-level modules. This dynamic also explains why the parameter ratio of {\textsc BLT}-Entropy ps=8 relative to the Llama 2 model increases modestly from 1.6x to 1.7x with model scaling.

To conclude, our investigation into patch-length scaling reveals that the {\textsc BLT} patch-based model exhibits superior scaling behavior when patch size and overall model capacity are increased together. These scaling benefits appear not only to endure but also to intensify as model sizes grow.

\section{Byte Modeling Improves Robustness} We further assess the robustness of the {\textsc BLT} architecture in contrast to token-based counterparts that do not leverage explicit byte-level data, and introduce a method for converting pretrained token-based models to process byte inputs.

\subsection{Character-Level Tasks}

One of the initial driving forces behind the development of byte-level models was their increased resilience to noise at the byte level within input data, as well as their capability to recognize the internal composition of tokens—a feature where existing tokenizer-based models often exhibit shortcomings. To assess these capabilities, we conduct supplementary experiments using benchmarks specifically designed to test both resilience against noise in the input and sensitivity to the makeup of individual characters in various languages, such as English and others, as well as numerical digits and phonemes. The outcomes of these experiments are summarized in Table~\ref{tab:char_tasks}.

\paragraph{Noisy Data} To assess the resilience of {\textsc BLT} compared to tokenizer-based architectures, we generate noisy variants of the benchmark classification datasets outlined in Section~\ref{section:evals}. Specifically, we introduce five separate character-level perturbations: (a) \textit{AntSpeak}: the text is rendered in uppercase with individual characters separated by spaces; (b) \textit{Drop}: 10\% of characters are randomly deleted; (c) \textit{RandomCase}: the casing of each character is randomized so that approximately half appear in uppercase and half in lowercase; (d) \textit{Repeat}: a random 20\% subset of characters is repeated up to four times; (e) \textit{UpperCase}: all text is converted to uppercase. Each of these noise types is independently applied to either the prompt, the completion, or both, and we report results as averages over these scenarios. Table \ref{tab:char_tasks} displays performance for the noised HellaSwag~\citep{Eval_hellaswag} task, showing that {\textsc BLT} consistently surpasses tokenizer-based models with respect to robustness, achieving an average margin of 8 points relative to a model trained on identical data, and even exceeding the results of the Llama 3.1 model, despite the latter's much larger training corpus.

\paragraph{Phonology - Grapheme-to-Phoneme (G2P)} We evaluate the effectiveness of {\textsc BLT} in converting a series of graphemes (alphabetic representations of words) into their corresponding phonemic transcriptions. The G2P task results, as shown in Table \ref{tab:char_tasks}, are obtained in a 5-shot experimental setup utilizing Phonology Bench~\citep{suvarna2024phonologybench}. Our findings indicate that {\textsc BLT} surpasses the performance of the Llama 3 1T tokenizer-based baseline on this particular task.

\paragraph{CUTE}
To evaluate character-level capabilities, we test {\textsc BLT} on the CUTE benchmark~\citep{edman2024cute}, which encompasses a range of tasks falling under three primary domains: compositional understanding, orthographic similarity, and sequence manipulation. These tasks are particularly demanding for models based on tokenization, which typically know the spellings of their tokens but lack the ability to leverage this knowledge effectively for text manipulation. As reported in Table~\ref{tab:char_tasks}, {\textsc BLT}-Entropy surpasses both BPE Llama 3 variants by over 25 points on this suite. Notably, our model achieves nearly perfect results (99.9\%) on the spelling-related tasks, highlighting its strength in character-level manipulations. The substantial gains are achieved despite {\textsc BLT} being trained on only one-sixteenth the data of Llama 3.1, suggesting that BPE-based architectures find it intrinsically challenging to acquire robust character-level representations. Figure~\ref{fig:cute_examples} presents examples where Llama 3, relying on a tokenizer, falters, whereas our {\textsc BLT} model succeeds. The only exceptions are word deletion and insertion tasks, where BPE-based models have an edge, likely because such operations are more naturally aligned with their structure; however, this performance difference is relatively minor. Moreover, progressing from characters to word-level understanding may be more tractable for byte-based models than vice versa. The evaluation procedure and prompts are consistent across models, utilizing the original Huggingface prompts. It should also be noted that BPE models could potentially see gains from further prompt engineering.

\paragraph{Low Resource Machine Translation} We assess the performance of {\textsc BLT} for translation tasks involving six prominent language families and twenty-one low-resource languages, each representing different scripts, utilizing the FLORES-101 benchmark~\citep{goyal2022flores}. SentencePiece BLEU scores are provided in Table~\ref{tab:flores}. The findings reveal that {\textsc BLT} achieves higher BLEU scores than a model employing the Llama 3 tokenizer, with a margin of 2 BLEU points for translations into English and 0.5 points for translations from English. For widely used language pairs, {\textsc BLT} shows results that are on par with, or marginally exceed, those of Llama 3. Notably, for numerous low-resource language pairs, {\textsc BLT} demonstrates superior performance, highlighting the advantages of byte-level modeling in addressing the challenges of less-represented byte sequences.

\subsection{Training {\textsc BLT} via Llama 3 Initialization}
\label{sec:distillation}
We investigate a methodology whereby {\textsc BLT} models capitalize on pre-existing, pre-trained models with tokenization to facilitate enhanced and more rapid convergence during training. This is accomplished by setting the initial global transformer parameters of {\textsc BLT} to match those from a pre-trained Llama 3.1 model. During the training process, the global transformer’s weights are fine-tuned with a learning rate that is one-tenth that used for both the local encoder and local decoder components, over the optimal number of steps as determined for Llama 3. Performance outcomes, detailed in Table~\ref{tab:distill}, show a clear advantage: {\textsc BLT} initialized from Llama 3.1 exhibits notable improvements relative to both the standard Llama 3 model and a baseline {\textsc BLT}, despite being allotted equivalent computational resources (measured in flops). Furthermore, in comparison to the {\textsc BLT}-Entropy model described in Table~\ref{tab:evals}—which was trained on a much larger dataset of 1T tokens—{\textsc BLT} initialized from Llama 3.1 still demonstrates better results on the MMLU benchmark. This finding indicates the effectiveness of this approach for reducing training compute requirements without sacrificing downstream performance.

This configuration can be interpreted as adapting models that utilize tokenizers into versions that do not require tokenizers, essentially transforming a pre-trained LLaMA 3.1 into a {\textsc BLT} model. For a thorough evaluation, we list the original LLaMA 3.1, which was trained on 15T tokens, in Table \ref{tab:distill} and compare its results with those of the {\textsc BLT} produced from LLaMA 3. The modified model shows a minor reduction in scores on MMLU and HumanEval, with more pronounced declines observed in other benchmarks. These outcomes indicate a need for additional research to more effectively exploit the pre-trained model and enhance its capabilities, especially by fine-tuning data mixtures and optimizing hyperparameters.

\section{Ablations and Discussion}
\label{section:ablations}
This section presents ablation studies that substantiate the rationale for the architectural decisions in {\textsc BLT}, alongside the patching approach and hyper-parameter selection for the {\textsc BLT} model with 8B parameters, which was trained on the {\textsc BLT}-1T dataset.

\paragraph{Entropy Model Hyper-parameters} To analyze how scaling the entropy model's size and varying the context window affect performance, we experiment with byte-level entropy transformer models across a parameter range from 1 million to 100 million, using context window lengths between 64 and 512. We display bits-per-byte (bpb) against training compute (flops) scaling laws for our $400m$ and $1b$ parameter {\textsc BLT} models trained on the Llama-2 dataset, as shown in \autoref{fig:entropy_ablation}. Our results indicate that both increasing model size and lengthening the context window enhance scaling performance; however, gains start to taper off past approximately 50 million parameters.

\paragraph{Types of Patching}

We conduct an ablation study of the four patching strategies outlined in Section~\ref{section:patching}: 1) Strided Patching using strides of 4 and 6, 2) Whitespace-based patching, 3) Patching via the Llama 3 BPE Tokenizer, and 4) Entropy-driven patching utilizing a lightweight byte LLM.

Although dynamic patching shortens the actual sequence length, we standardize the sequence length to ensure that each patching scheme maintains a comparable context window. During both training and inference, all models process an equivalent expected number of bytes per sequence, which controls for confounding variables related to handling extended contexts. The outcomes of these ablation studies are depicted in Figure~\ref{fig:scaling-compute-opt}. Notably, every alternative patching strategy surpasses static patching in performance, with space patching closely matching the effectiveness of dynamic entropy based patching.

In \autoref{tab:evals_ablation}, we report benchmark results for {\textsc BLT} models, examining the performance of tokenizer-based approaches, space patching, and entropy-based patching, each trained on the Llama 2 dataset for an empirically determined number of training steps~\citep{dubey2024llama}. While space patching offers a more straightforward technique by not requiring real-time entropy modeling during training, our results indicate that the advantages seen with entropy-based patching in the context of scaling behavior~(Section~\ref{section:scaling}) persist and extend to downstream evaluation benchmarks as well.\footnote{Space patching outcomes originate from prior experiments that excluded cross-attention; however, consistent patterns are observed when cross-attention is incorporated.}

\paragraph{Cross-Attention}
\autoref{tab:crossattn_ablations_1b} presents an ablation study on integrating cross-attention modules at different locations within the encoder and decoder of {\textsc BLT}. Regarding encoder cross-attention, we explore three initialization strategies for the queries: (1) assigning every global state an identical learned embedding, (2) utilizing a hash embedding based on the patch's byte content, and (3) employing a pooled representation of the patch bytes derived from the encoder's hidden states at the corresponding layer.

Our results indicate that incorporating cross-attention within the \textit{decoder} yields the best performance. Employing cross-attention in the encoder offers marginal gains, but these occur only when queries are initialized via pooling. Furthermore, we observe that cross-attention provides notable benefits on the Common-Crawl dataset, especially when larger patch sizes are used.

\paragraph{n-gram Hash Embeddings} 

We conduct ablation studies using n-gram hash embedding vocabularies of sizes 0, 100K, 200K, and 400K, with the results summarized in Table~\ref{tab:ngram_ablations_1b}. Our analysis shows that incorporating hash embeddings consistently benefits all tested domains, with notably larger gains observed for Wikipedia and Github (yielding a 0.04 bpb improvement compared to a 0.01 bpb gain after 15k steps at the 8B parameter setting). When scaling to 8B parameters, reducing the number of hashes from 500K to 300K results in a marginal performance decrease of approximately 0.001 bpb after 15k training steps. This suggests that hash embeddings are essential for elevating the performance of {\textsc BLT} to the level of models using tokenizers. However, increasing the hash vocabulary beyond 300K yields progressively smaller benefits. Moreover, our results indicate that improvements from hash embeddings largely supplement those from cross-attention, as each contributes advantages to distinct datasets.

\paragraph{Local Model Hyperparameters} Table \ref{tab:local_layers} presents an ablation of different configurations for the number of layers used in both the local encoder and decoder. Notably, when employing hash n-gram embeddings, {\textsc BLT} demonstrates strong performance even when the encoder is minimal—utilizing only a single layer—while the decoder benefits from a greater number of layers.

\section{Related Work}
\label{section:related_work}

\textbf{Character-Level RNNs:} Character Language Modeling has maintained its significance in neural modeling research~\citep{sutskever2011generating,mikolov2012subword,graves2013generating}, primarily because it naturally handles out-of-vocabulary terms without relying on fallback mechanisms. \cite{kim2016character} introduce a model employing convolutions and highway networks over input characters, which are then fed into LSTM-based RNNs. Their approach not only matches the then-leading RNN language models' performance on English but also surpasses them in morphologically complex languages, highlighting another key strength of character-level LLMs. Similarly, \cite{kenter2018byte} leverage byte-level LSTM architectures for machine comprehension and demonstrate superior results over word-based systems on languages such as Turkish and Russian, both of which have intricate morphology. In a related effort, \cite{zhang2015character} apply character-based convolutional networks to text classification tasks and show that, for certain problems, these outperform traditional word-based counterparts. Further enhancing these approaches, \citet{chung2019hierarchical} develop hierarchical LSTM architectures with boundary detectors at each hierarchical stage to uncover the latent structure within text, which yields improved character-level language modeling outcomes. In contrast to attention-based methods, ByteNet by \cite{kalchbrenner2016neural} incorporates CNN layers over character sequences for machine translation tasks.

\textbf{Character-Level Transformers:} The introduction of attention-based transformer architectures~\citep{vaswani2017attention}, in conjunction with subword tokenization approaches~\citep{sennrich-etal-2016-neural}, led to substantial advances in neural language models across a range of language modeling benchmarks. Nevertheless, the utilization of word and sub-word units introduces a particular inductive bias, constraining the abstraction level at which models learn. Seeking to integrate the achievements of transformer models with encouraging early outcomes in character-level modeling, \citet{al2019character} employed extremely deep transformers augmented by auxiliary loss functions, resulting in transformer-based architectures that surpassed prior LSTM-based character language models. Despite this progress, there remained a notable performance disparity when compared to word-level large language models. Similarly, GPT-2~\citep{radfordgpt2} reported that byte-level language models, even when trained on extensive datasets such as the 1 Billion Word Benchmark, did not reach the competitive effectiveness of their word-level counterparts.

\cite{Choe2019BridgingTG} showed that byte-level LLMs built upon transformer architectures can surpass subword-level models with a similar number of parameters; however, these models demand substantially greater computational resources and require significantly longer training times. In a related vein, \cite{el2020characterbert} introduced CharFormer, a BERT-based model that leverages convolutions over character embeddings to construct word-level representations, reporting enhanced performance in the medical domain—albeit with considerably higher computational costs. In another advancement, \cite{clark2022canine} presented CANINE, an encoder-only model with 150M parameters designed to process raw character sequences. The CANINE model integrates a deep transformer core, which aligns conceptually with the architecture of our global model, together with a combination of localized transformers and strided convolutions for input character downsampling. This architecture allows CANINE to exceed the performance of token-level encoder-only baselines like mBERT on downstream multilingual evaluation benchmarks. 

ByT5~\citep{xue2022byt5} investigated byte-level encoder-decoder models without employing any patching mechanisms. Their results demonstrated enhanced robustness to input noise and competitive performance relative to tokenizer-dependent models, even when trained with four times less data. Nonetheless, the absence of patching operations resulted in the model having to carry out costly attention computations over every input byte, which led to prohibitive computational overhead. Shifting from subword units to byte-level modeling inherently expands input sequence lengths, presenting significant challenges for the scalable and efficient training of such models.

More recently, the introduction of the Mamba Architecture~\citep{gu2023mamba}, which features the ability to preserve a fixed-size memory state across extremely long contexts, has led to new directions. \cite{wang2024mambabyte} utilized this architecture to develop a byte-level Mamba-based model—again, eschewing patching techniques—and demonstrated that it can surpass byte-level transformer models with controlled compute budgets at the 350M parameter scale, as measured by bits-per-byte across multiple datasets.

\textbf{Patching-based approaches:} Patching methods have been shown to reduce the otherwise high computational requirements (flops) associated with byte-level LLMs, without necessarily compromising their effectiveness. Several studies have reported promising outcomes with such techniques on relatively modest model sizes and limited training data. For instance, \cite{nawrot-etal-2022-hierarchical} introduced static patching involving both downsampling and upsampling operations, proposing the hourglass transformer architecture, which surpasses other byte-level benchmarks at the 150M parameter level. Building on this, \cite{nawrot-etal-2023-efficient} advanced the methodology through dynamic patching, introducing innovations such as a boundary-predictor learned end-to-end, variants of the predictor supervised with select tokenizers, and an entropy-based patching strategy akin to {\textsc BLT}. Their results indicate that these dynamic schemes outperform contemporaneous vanilla transformers in language modeling at the 40M parameter scale with roughly 400M training tokens. Additionally, \cite{lester2024training} explored the effects of compressing sequences via arithmetic coding to attain compression rates exceeding those of BPE, and through the adoption of an equal-info windows method, managed to improve over byte-level baselines in language modeling tasks, though their approach did not surpass subword-based baselines.

This study takes significant cues from MegaByte~\citep{yu2023megabyte}, a causal LLM relying solely on a decoder architecture. MegaByte utilizes a predetermined, static approach for dividing inputs into patches and concatenating their representations, facilitating the conversion from bytes to patches. On the decoding side, a localized model is responsible for translating patches back into their original byte sequences. The authors of MegaByte have shown that, at the 1B parameter level and over datasets encompassing approximately 400B bytes, this framework can perform on par with models utilizing tokenization.

Within all our experimental analyses, we systematically examine MegaByte, noting that the use of fixed patching trails behind state-of-the-art tokenizer-based models that have been optimally trained for a given computation budget. Our results highlight how {\textsc BLT} narrows this performance disparity. Similar conclusions are reached by \cite{slagle2024spacebyte}, who advocate for adapting the static patching approach to operate over whitespace and other space-related bytes, along with incorporating a localized encoder component. Their findings reveal computationally controlled performance gains over token-based transformers in specific contexts, such as code (Github) and scientific literature (arXiv) at the 1B parameter level. We extend our evaluation to include this model and provide evidence that more substantial architectural enhancements are required if byte-level approaches are to scale and rival leading token-based models, exemplified by Llama 3.

\section{Limitations and Future Work}

In this study, model training is conducted for the empirically optimal number of steps identified in Llama~3~\citep{dubey2024llama}, to inform architectural decisions. Nonetheless, these scaling laws were originally derived for BPE-level transformer models, and may not yield ideal ratios of data to parameters when applied to {\textsc BLT}. Deriving scaling laws specific to {\textsc BLT} is a topic reserved for future investigation, and doing so may reveal even more advantageous scaling dynamics for our approach. Furthermore, our current experimental scope is primarily limited to model sizes up to 1B parameters; it is plausible that the most effective architectural configurations could shift as we scale models to 8B parameters or larger, possibly resulting in enhanced performance at these higher capacities.

Current transformer frameworks and repositories prioritize optimization for architectures utilizing tokenizers. Although our work provides theoretical flop-matched comparisons and leverages some optimized implementations—such as FlexAttention—for layers diverging from the standard transformer configuration, our practical implementations might still lag behind tokenizer-based models in terms of actual wall-clock performance and could be improved with additional optimization efforts.

Although {\textsc BLT} currently relies on a separately pre-trained entropy model for patching, a promising avenue for future research lies in integrating the learning of the patching model into an end-to-end training framework. Section \ref{sec:distillation} provides preliminary evidence that supports the feasibility of “byte-ifying” models based on tokenizers, such as Llama 3—which is trained on datasets exceeding 10T tokens—by transferring and freezing the pretrained weights of the global transformer component. Advancing this approach could potentially reveal techniques that preserve the advantages of bytefying while also surpassing the performance of existing tokenizer-based models, all without the need for retraining from the ground up.

\section{Conclusion}
\label{section:conclusion}

In this work, we introduce the Byte Latent Transformer (\textbf{BLT}), an innovative architecture that moves beyond the standard reliance on fixed-vocabulary tokenization found in most large language models. By leveraging a flexible, trainable mechanism for assembling bytes into variable-length patches, {\textsc BLT} dynamically assigns computational resources in accordance with data complexity, thereby substantially enhancing both computational efficiency and model robustness. Through comprehensive scaling experiments, we show that {\textsc BLT} models rival the performance of established token-based architectures such as Llama 3 when scaled to 8B parameters and 4T bytes of training data, while offering up to a 50\% reduction in inference flops at only a marginal cost to evaluation scores. Notably, {\textsc BLT} introduces a new axis for efficient scaling by permitting simultaneous enlargement of both the model and the patch size within a constant inference compute budget—a property well-suited to realistic compute constraints. By operating directly on byte-level representations, {\textsc BLT} not only enhances the model’s robustness to noise and its capacity to process rare or atypical inputs, but also fosters a richer grasp of underlying sub-word features. Collectively, these findings establish {\textsc BLT} as a compelling alternative to traditional tokenization-based strategies, yielding a flexible, resilient, and efficient foundation for future language models.

\end{document}